\clearpage\SourceLeaf{85}%
\pdfbookmark[1]{Berechnung des Pessachfestes}{osterfest1802}
\begin{center}
{\large BERECHNUNG DES PESSACHFESTES.\par}
\medskip\hrule\medskip
{\small Monatliche Correspondenz zur Beförderung der Erd- und Himmels-Kunde,\par
herausgegeben vom Freiherrn \textsc{von Zach}. Mai 1802.\par}
\medskip\hrule
\end{center}\medskip
Der 15. Nisan des jüdischen Jahres $A$, an dem die Juden ihr Pessachfest feiern, fällt in das Jahr $A-3760=B$ der christlichen Zeitrechnung. Zur Bestimmung des entsprechenden Monatstages dient die folgende rein arithmetische Regel.

Man teilt $12A+17$ oder, was hier dasselbe ergibt, $12B+12$ durch 19 und nennt den Rest $a$. Dann teilt man $A$ oder $B$ durch 4 und setzt den Rest $=b$. Man berechnet den Wert
\begin{center}\small
\begin{tabular}{lll}
\multicolumn{1}{c}{von} & \multicolumn{1}{c}{oder von} & \shortstack{Wert der Dezimalbrüche\\in gewöhnlichen Brüchen}\\[5pt]
$32.0440932$ & $20.0955877$ & $32\frac{4343}{98496}\quad20\frac{9415}{98496}$\\[5pt]
$+1.5542418a$ & $+1.5542418a$ & $1\frac{272953}{492480}$\\[5pt]
$+0.25\quad b$ & $+0.25\quad b$ & $\frac14$\\[5pt]
$-0.003177794A$ & $-0.003177794B$ & $\frac{313}{98496}$
\end{tabular}
\end{center}
und setzt ihn $=M+m$, sodass $M$ den ganzzahligen Anteil und $m$ den (dezimalen) Bruchteil bezeichnet. Schließlich teilt man $M+3A+5b+5$ oder $M+3B+5b+1$ durch 7 und setzt den Rest $=c$. Nun sind die folgenden vier Fälle zu unterscheiden.

I. Ist $c=2$ oder beträgt der Wert 4 oder 6, so fällt Pessach auf den $M+1^{\mathrm{ten}}$ März \textit{alten Stils}. Dafür schreibt man den $M-30^{\mathrm{ten}}$ April, wenn wegen \textit{Adu} die Bedingung $M>30$ erfüllt ist.
\SourceLeaf{86}%

II. Ist $c=1$, zugleich $a>6$ und außerdem $m\geqq0.63287037$
$(\frac{311676}{492480}=\frac{1367}{2160})$, so fällt Pessach wegen \textit{Getred} auf den $M+2^{\mathrm{ten}}$ März \textit{a. St.}

III. Ist $c=0$, zugleich $a>11$ und außerdem $m\geqq0.89772376$\hfill\break
\mbox{$(\frac{442111}{492480}=\frac{23269}{25920})$}, so fällt Pessach wegen \textit{Batu Thakpad} auf den $M+1^{\mathrm{ten}}$ März \textit{a. St.}

IV. In allen übrigen Fällen fällt Pessach auf den $M^{\mathrm{ten}}$ März alten Stils.

\begin{center}\rule{.25\linewidth}{.3pt}\end{center}

\textit{Erste Anmerkung.} Diese Regeln dienen zugleich zur Bestimmung des 1. \textit{Tisri}, also des Neujahrstags, der immer 163 Tage nach Pessach des vorhergehenden Jahres liegt.

\textit{Zweite Anmerkung.} Das Jahr $A$ ist ein gewöhnliches Jahr (mit 12 Monaten), wenn $a<12$, und ein Schaltjahr (mit 13 Monaten), wenn $a>11$.

\begin{center}\small
\begin{tabular}{lll}
\textit{Beispiel zu diesen Regeln:} & $A=5562$ & $B=1802$\\
& $12A+17=66761$ & $12B+12=21636$\\
\multicolumn{3}{c}{Die Division durch 19 ergibt $a=14$}\\
\multicolumn{3}{c}{Die Division von 5562 oder 1802 durch 4 ergibt $b=2$.}\\
Damit liefert die obige Formel & $32.0440932$ & $20.0955877$\\
& $+21.7593852$ & $21.7593852$\\
& $+\phantom{0}0.5$ & $0.5$\\
& $-17.6748903$ & $-\phantom{0}5.7263848$\\\cline{2-3}
& $36.6285881$ & $36.6285881$
\end{tabular}
\end{center}
also $M=36,\quad m=0.6285881$,
\[
\begin{aligned}
M+3A+5b+5&=36+16686+10+5=16737\\
M+3B+5b+1&=36+\phantom{0}5406+10+1=\phantom{0}5453
\end{aligned}
\]
und die Division durch 7 ergibt $c=0$.

Da $m$ hier kleiner als $0.89772376$ ist, greift Regel III hier nicht. Nach Regel IV fällt Pessach daher auf den 36. März alten Stils oder den 48. März neuen Stils, das heißt auf den 17. April.

In den \textit{meisten} Fällen genügt es, mit der obigen Formel nur etwa 2 Dezimalstellen zu berechnen.

\begin{center}\rule{.25\linewidth}{.3pt}\end{center}
